\documentclass[11pt]{article}
\usepackage[margin=1in]{geometry}
\usepackage{amsmath,amssymb,amsthm,booktabs,xurl,hyperref}
\hypersetup{colorlinks=true,linkcolor=blue,urlcolor=blue,
 pdftitle={An exact exponent obstruction for lower-end projection-cone simplex stability},
 pdfauthor={Research supplement to entry005}}
\newtheorem{theorem}{Theorem}
\newtheorem{corollary}[theorem]{Corollary}
\newtheorem{proposition}[theorem]{Proposition}
\newcommand{\conv}{\operatorname{conv}}
\newcommand{\BM}{d_{\mathrm{BM}}}
\newcommand{\vol}{\operatorname{vol}}
\newcommand{\ee}{\mathfrak e}
\title{An exact exponent obstruction for lower-end projection-cone simplex stability}
\author{Research supplement to entry005; AI-assisted research note}
\date{7 October 2026}
\begin{document}
\maketitle

\noindent\textbf{Outcome.} In every fixed dimension $d\ge3$, any power estimate for
the lower-end defect of entry005, whether in Banach--Mazur distance to any simplex
or in centroid containment by every maximum-volume inscribed simplex, must have
exponent at most $1/(d-1)$. The example also gives an exact lower bound on the
constant in the latter estimate at this exponent. This is a necessary limit,
not a claim that the endpoint exponent has been proved for arbitrary bodies.
The existing explicit power theorem in the quantitative supplement is not reproved here.

\section{Conventions and public-source scope}

For a full-dimensional convex body in $\mathbb R^d$, use exactly the public
entry005 convention
\[
 R(K)=\frac{|\Pi K|}{|K|^{d-1}},\qquad
 a(K)=\left(\frac d{d+1}\right)^d\frac{R(\mathcal P K)}{R(K)}-1,
 \qquad \ee(K)=a(K)-\frac1{d+1}.
\]
Here $\mathcal P K$ is a pyramid over $K$ and $\Pi K$ is its projection body.
The definition and the determinant formula below were read in public v3;
v4 preserves the definition. The latest public v5 at commit
\nolinkurl{31e3d8a37e4a3051a7f5a2535be1c75642e15bcc} concerns balanced
product/join recursions and does not replace the lower-end definition or add
a quantitative lower-end theorem within that versioned package.

The additive quantitative supplement is available publicly
at commit \nolinkurl{6785c1c830f8e19e2eb07b0bb89f4d475a8b154a}.
Its source defines $a$ and $E$ exactly as above and proves the
every-maximum-simplex bound with exponent
$1/[d(d^2+2d+2)]$. Its stated inputs are the v2--v3 invariant identities
and classical Cauchy and mixed-volume formulas; its manifest pins the
audited dependency commit
\nolinkurl{dd969ea280508a15995fbe83a80cac0f6d227f73}.
This note does not repeat that quantitative proof.

For a simplex $S\subset K$ with centroid $g_S$, define the \emph{excess}
centroid dilation by
\[
 E(K,S)=\inf\{\epsilon\ge0:K\subset g_S+(1+\epsilon)(S-g_S)\}.
\]
Thus the dilation factor itself is $1+E$. Write $M(K)$ for the maximum volume
of an inscribed $d$-simplex. For any fixed reference simplex $\Delta_d$, the
usual affine Banach--Mazur distance to its affine class can equivalently be
written
\[
 \BM(K,\Delta_d)=\inf\{\lambda\ge1:
 S\subset K\subset v+\lambda S,
 \ S\text{ a full-dimensional simplex},\ v\in\mathbb R^d\}.
 \tag{1}
\]
The free translation in (1) and optimization over all $S$ are important.
They are both absent from the prescribed-centroid, prescribed-maximum-simplex
quantity $E$.

For a polytope with facet area times outward unit normal $u_i$ and facet area
times support number $b_i$, public v3 gives
\[
 a(K)=\frac{\displaystyle\sum_{|I|=d+1}
       |\det((u_i,b_i):i\in I)|}
 {d|K|\displaystyle\sum_{|J|=d}|\det(u_j:j\in J)|}.
 \tag{2}
\]
It is legitimate to use (2) before translating an interior point to zero:
a translation by $z$ replaces $b_i$ by $b_i-u_i\cdot z$, an elementary row
operation in every lifted determinant. No nonpositive $b_i$ is used as a
probability weight in this note.

\section{The obstruction family and its complete calculation}

Let $e_1,\ldots,e_d$ be the coordinate vectors and, for $0<t<1$, set
\[
 K_t=\{x\in\mathbb R^d:x_i\ge0,\ t\le x_1+\cdots+x_d\le1\}
     =\conv\{e_1,\ldots,e_d,te_1,\ldots,te_d\}.
 \tag{3}
\]
This is a standard simplex with a homothetic cap cut from one vertex.

\begin{theorem}\label{thm:exact}
For every integer $d\ge2$ and $0<t<1$, the following assertions hold.

\medskip\noindent
\textup{(i)} With $q=t^{d-1}$, the invariant is exactly
\[
 a(K_t)=
 \frac{1+(d-1)q-(d-1)t^d-qt^d}
 {(1-t^d)(d+1+(d-1)q)},
 \tag{4}
\]
and consequently
\[
 \ee(K_t)=
 \frac{t^{d-1}\bigl[d(d-1)-(d+1)(d-2)t-2t^d\bigr]}
 {(d+1)(1-t^d)(d+1+(d-1)t^{d-1})}.
 \tag{5}
\]
In particular,
\[
 \lim_{t\downarrow0}\frac{\ee(K_t)}{t^{d-1}}
   =c_d:=\frac{d(d-1)}{(d+1)^2}>0.
 \tag{6}
\]

\medskip\noindent
\textup{(ii)} The maximum simplex volume is
\[
 M(K_t)=\frac{1-t}{d!}.
 \tag{7}
\]
Every maximum-volume inscribed simplex, including those with vertices in the
relative interiors of faces, has the form
\[
 S_p=\conv\{e_1,\ldots,e_d,tp\},\qquad
 p_i\ge0,\quad\sum_{i=1}^d p_i=1.
 \tag{8}
\]
For each such simplex,
\[
 E(K_t,S_p)=(d+1)t\max_i p_i.
 \tag{9}
\]
Thus
\[
 \inf_{S\text{ maximum}}E(K_t,S)=\frac{d+1}{d}t,
 \qquad
 \sup_{S\text{ maximum}}E(K_t,S)=(d+1)t.
 \tag{10}
\]

\medskip\noindent
\textup{(iii)} The distance to any simplex satisfies
\[
 \left(1+t+\cdots+t^{d-1}\right)^{1/d}
 \le \BM(K_t,\Delta_d)\le 1+t.
 \tag{11}
\]
In fact, for every fixed maximum simplex $S_p$, its smallest enclosing
\emph{translated} homothet has factor exactly $1+t$.
In particular, for $0<t<1$,
\[
 \frac{t}{2d}\le \BM(K_t,\Delta_d)-1\le t,
 \qquad
 \liminf_{t\downarrow0}\frac{\BM(K_t,\Delta_d)-1}{t}\ge\frac1d.
 \tag{12}
\]

\medskip\noindent
\textup{(iv)} The Minkowski asymmetry is exactly
\[
 s(K_t):=\inf\{r\ge1:\exists z,
              -(K_t-z)\subset r(K_t-z)\}=d-t.
 \tag{13}
\]
\end{theorem}

\begin{proof}
We prove each assertion directly.

\medskip\noindent\textbf{Facet determinants.}
Put $c=1/(d-1)!$ and $r=1-q$. The $d$ coordinate facets have area
$c(1-t^{d-1})$. The top and bottom facets have areas $c\sqrt d$ and
$c\sqrt d\,t^{d-1}$, respectively. The facet data in the coordinates (3) are
therefore
\[
 (u_i,b_i)=(-cr e_i,0)\quad(1\le i\le d),\qquad
 (u_+,b_+)=(c\mathbf1,c),\qquad
 (u_-,b_-)=(-cq\mathbf1,-ct^d).
 \tag{14}
\]
The volume is $|K_t|=(1-t^d)/d!$, so $d|K_t|=c(1-t^d)$.

For horizontal $d$-minors, the all-coordinate subset contributes $c^d r^d$;
the $d$ subsets with the top facet contribute $c^d r^{d-1}$ each; and the
$d$ subsets with the bottom facet contribute $c^d q r^{d-1}$ each.
Subsets with both top and bottom have determinant zero because their normals
are parallel. Hence the horizontal sum is
\[
 H=c^d r^{d-1}\bigl[d+1+(d-1)q\bigr].
 \tag{15}
\]
For the lifted $(d+1)$-minors, the subset containing all coordinate facets
and the top facet contributes $c^{d+1}r^d$. Replacing the top facet by the
bottom facet gives $c^{d+1}r^dt^d$. Each of the $d$ subsets with top, bottom and $d-1$
coordinate facets contributes
$c^{d+1}r^{d-1}(q-t^d)$; here $q-t^d=q(1-t)>0$.
Thus the lifted sum is
\[
 L=c^{d+1}r^{d-1}
   \bigl[1+(d-1)q-(d-1)t^d-qt^d\bigr].
 \tag{16}
\]
Dividing $L$ by $d|K_t|H$ gives (4); subtracting $1/(d+1)$ gives (5).
The bracket in (5) decreases strictly to zero as $t\uparrow1$, so the defect
is positive. Dividing (5) by $t^{d-1}$ and passing to $t=0$ proves the exact
limit (6), without any numerical fitting.

\medskip\noindent\textbf{Maximum simplices, including nonvertex ones.}
First consider a simplex selected from the $2d$ vertices in (3). If a ray
$\mathbb R_+e_i$ is unrepresented, the simplex lies in the coordinate
hyperplane $x_i=0$ and has zero volume. Otherwise $d+1$ distinct vertices
on $d$ rays give exactly one doubled ray, with vertices $e_j,te_j$, and one
vertex on every other ray. Subtract the two doubled columns in the lifted
vertex determinant. If $k$ of the other $d-1$ selected vertices are bottom
vertices, its absolute determinant is
\[
 (1-t)t^k.
 \tag{17}
\]
The largest value is $1-t$, attained precisely by all $d$ top vertices and
one bottom vertex. Repeated vertices give determinant zero.

Now let $z_0,\ldots,z_d$ be arbitrary vertices of an inscribed simplex.
Represent each $z_j$ as a convex combination of the $2d$ vertices, and
independently sample a vertex $Y_j$ from that combination. Multiaffinity of
the lifted determinant $D$ gives
\[
 |D(z_0,\ldots,z_d)|
 =|\mathbb E D(Y_0,\ldots,Y_d)|
 \le\mathbb E|D(Y_0,\ldots,Y_d)|\le1-t.
 \tag{18}
\]
This proves the upper bound (7), already attained by the preceding examples.
If the simplex is maximum, equality in (18) implies that
$|D(Y_0,\ldots,Y_d)|=1-t$ almost surely. By (17), exactly one of the $Y_j$
is a bottom vertex and the other $d$ are the $d$ distinct top vertices,
almost surely.

Let $I_j$ be the indicator that $Y_j$ is a bottom vertex. Independence and
$\sum_j I_j=1$ almost surely imply
\[
 0=\operatorname{Var}\Bigl(\sum_j I_j\Bigr)
  =\sum_j \mathbb P(I_j=1)\bigl[1-\mathbb P(I_j=1)\bigr].
\]
Thus every slot has deterministic top/bottom type, with precisely one bottom
slot. The supports of the $d$ top slots must be pairwise disjoint; otherwise
independent sampling gives a repeated top vertex with positive probability.
There are $d$ nonempty disjoint supports in a set of $d$ top vertices, so
each is a singleton. The bottom slot is any convex combination $tp$.
This proves (8). Conversely, the height above the hyperplane
$\sum x_i=1$ of any point $tp$ is fixed, so every simplex (8) has volume
$(1-t)/d!$.

\medskip\noindent\textbf{Centroid containment.}
Write $s_x=\sum_i x_i$. Barycentric coordinates for $S_p$, with index $0$
for $tp$, are
\[
 \beta_0(x)=\frac{1-s_x}{1-t},\qquad
 \beta_i(x)=x_i-tp_i\beta_0(x)\quad(1\le i\le d).
 \tag{19}
\]
Evaluation at the $2d$ vertices of $K_t$ shows
\[
 \min_{K_t}\beta_0=0,\qquad \min_{K_t}\beta_i=-tp_i.
 \tag{20}
\]
For the latter minimum, use $te_j$ with $j\ne i$, available since $d\ge2$.
Affine functions attain their extrema on these vertices.
The centroid-dilated simplex $g_{S_p}+(1+\epsilon)(S_p-g_{S_p})$
is exactly the region $\beta_i\ge-\epsilon/(d+1)$ for every
$0\le i\le d$. Equations (20) give (9), and optimizing $\max_i p_i$
on the probability simplex gives (10).

\medskip\noindent\textbf{Translated containment and distance to any simplex.}
For barycentric coordinates of any simplex $S$, a translated homothet
$v+\lambda S$ has inequalities $\beta_i\ge a_i$ with
$\sum_i a_i=1-\lambda$. Conversely any such thresholds describe a
translated homothet. To contain $K$, one needs
$a_i\le\min_K\beta_i$ for each $i$, so
\[
 \lambda\ge1-\sum_i\min_K\beta_i.
 \tag{21}
\]
Choosing all thresholds equal to their minima realizes equality whenever
$S\subset K$. For (20), the sum is $-t$, giving factor exactly $1+t$.
Concretely,
\[
 K_t\subset -tp+(1+t)S_p.
 \tag{22}
\]
This proves the upper bound in (11) under the full freedom in (1).

For the lower bound, any simplex $S\subset K_t$ has $|S|\le M(K_t)$.
If $K_t\subset v+\lambda S$, volume monotonicity and (7) imply
\[
 \lambda^d\ge\frac{|K_t|}{|S|}
 \ge\frac{|K_t|}{M(K_t)}
 =\frac{1-t^d}{1-t}=1+t+\cdots+t^{d-1}.
 \tag{23}
\]
Taking the infimum proves the lower bound in (11). Furthermore
\[
 (1+t)^{1/d}-1
 =\frac1d\int_0^t(1+u)^{-(d-1)/d}\,du
 \ge\frac{t}{d\,2^{(d-1)/d}}\ge\frac{t}{2d}.
\]
Finally the derivative at zero of
$(1+t+\cdots+t^{d-1})^{1/d}$ is $1/d$, proving (12).
This lower bound concerns \emph{all} simplices, not only maximum simplices.
We do not assert an exact formula for the unrestricted distance in (11).

\medskip\noindent\textbf{Minkowski asymmetry.}
The containment in (13) is equivalent to
$-K_t+(r+1)z\subset rK_t$. Writing $Z=\sum_i z_i$, its necessary and
sufficient facet inequalities are
\[
 (r+1)z_i\ge1\quad(1\le i\le d),\qquad
 1+rt\le(r+1)Z\le r+t.
 \tag{24}
\]
Summing the coordinate inequalities and using the last upper bound gives
$r\ge d-t$. This value is attainable: take
$z_i=1/(d+1-t)$. The bottom inequality follows from
\[
 d-1-(d-t)t=(1-t)(d-1-t)\ge0.
\]
All other inequalities in (24) are equalities. This proves (13).
\end{proof}

\section{Universal implications and necessary constants}

\begin{corollary}\label{cor:power}
Fix $d\ge3$. A power estimate valid for every full-dimensional body of
dimension $d$, even restricted to a sufficiently small lower-end defect,
can have exponent $\alpha$ only if $\alpha\le1/(d-1)$, for each of:
\[
 \BM(K,\Delta_d)-1\le C_d\ee(K)^\alpha;
 \tag{25}
\]
\[
 \inf_{S\text{ maximum}}E(K,S)\le C_d\ee(K)^\alpha;
 \tag{26}
\]
\[
 E(K,S)\le G_d\ee(K)^\alpha
 \quad\text{for every maximum-volume inscribed simplex }S.
 \tag{27}
\]
At the critical exponent $1/(d-1)$, any constant in (27) must satisfy
\[
 G_d\ge(d+1)\left(\frac{(d+1)^2}{d(d-1)}\right)^{1/(d-1)}>d+1.
 \tag{28}
\]
Any constant in (26) must be at least $1/d$ times the right-hand side of
(28), and any constant in (25) must be at least
\[
 \frac1d\left(\frac{(d+1)^2}{d(d-1)}\right)^{1/(d-1)}.
 \tag{29}
\]
These are necessary constants, not claimed optimal universal constants.
In particular a dimension-independent constant is impossible in (27)
at the critical exponent.
\end{corollary}

\begin{proof}
Equations (6), (10), and (12) give a positive constant times $t$ on the
left of each estimate and a constant times $t^{(d-1)\alpha}$ on its right.
If $(d-1)\alpha>1$, division by $t$ and passage to $t\downarrow0$ is a
contradiction. At equality of exponents, take $p=e_1$ in (9) for (28),
$p_i=1/d$ for (26), and the lower-limit bound in (12) for (29).
\end{proof}

Combining the released upper estimate with Corollary \ref{cor:power} places
the supremum of admissible powers, for either the every-maximum-simplex
quantity or the unrestricted simplex Banach--Mazur distance, in the interval
\[
 \frac1{d(d^2+2d+2)}\ \le\ \alpha_d^*\ \le\ \frac1{d-1}.
 \tag{29a}
\]
For $d=3$ this is $1/51\le\alpha_3^*\le1/2$. The gap between these
two endpoints remains unresolved here.

\noindent\textbf{The tetrahedral case.}
For $d=3$, cancellation of $1-t$ in (5) gives the especially simple formula
\[
 \boxed{\ee(K_t)=
 \frac{t^2(t^2+t+3)}{4(1+t+t^2)(2+t^2)},\qquad
 \lim_{t\downarrow0}\frac{\ee(K_t)}{t^2}=\frac38.}
 \tag{30}
\]
Every maximum tetrahedron is $S_p=\conv\{e_1,e_2,e_3,tp\}$ and has
$E=4t\max p_i$. Thus no exponent greater than $1/2$ is possible;
at exponent $1/2$ the every-maximum-tetrahedron constant is at least
$4\sqrt{8/3}$, and the best-maximum-tetrahedron constant is at least
$(4/3)\sqrt{8/3}$.

\section{A precise obstruction to a direct Rogers--Shephard bridge}

\begin{proposition}\label{prop:rs}
For the same family in dimension three, put
\[
 \delta_{\rm RS}(K)=1-\frac{|K-K|}{20|K|}.
\]
Then
\[
 |K_t-K_t|=\frac{10-3t-6t^2-t^3}{3},\qquad
 \delta_{\rm RS}(K_t)=\frac{3t(1+3t)}{10(1+t+t^2)}.
 \tag{31}
\]
In particular,
\[
 \lim_{t\downarrow0}\frac{\delta_{\rm RS}(K_t)}t=\frac3{10},
 \qquad 3-s(K_t)=t.
 \tag{32}
\]
Neither $\delta_{\rm RS}(K)\le C\ee(K)$ nor
$3-s(K)\le C\ee(K)$ can hold uniformly near the lower endpoint.
More generally either of these two bridges with $\ee^\beta$ requires
$\beta\le1/2$.
\end{proposition}

\begin{proof}
For $x\in\mathbb R^3$, let $P=\sum_i(x_i)_+$ and $N=\sum_i(-x_i)_+$.
Every representation $x=a-b$ with $a,b\ge0$ has
$a=x_++w$, $b=x_-+w$ for some $w\ge0$.
Writing $W=\sum_i w_i$, membership $a,b\in K_t$ is equivalent to
\[
 \max(0,t-\min(P,N))\le W\le1-\max(P,N).
\]
Hence
\[
 K_t-K_t=\{P\le1,\ N\le1,\ |P-N|\le1-t\}.
 \tag{33}
\]
The difference body of the uncut simplex is $\{P\le1,N\le1\}$.
Its two pure-sign orthants each have volume $1/6$, and each of the six
mixed orthants has volume $1/2$. Its total volume is therefore $10/3$.
The two caps $P-N>1-t$ and $N-P>1-t$ are disjoint and have equal volumes.

For the first cap, the all-positive orthant contributes
$[1-(1-t)^3]/6$. Each of the three orthants with two positive coordinates
and one negative coordinate contributes
\[
 \int_{1-t}^1 P(P-1+t)\,dP=
 \frac{(1-t)t^2}{2}+\frac{t^3}{3}.
\]
Each of the three orthants with one positive coordinate and two negative
coordinates contributes
\[
 \frac12\int_{1-t}^1(P-1+t)^2\,dP=\frac{t^3}{6}.
\]
These integrals follow by summing the nonnegative coordinates: two such
coordinates have slice density equal to their sum.
The cap volume is consequently $t/2+t^2+t^3/6$. Subtracting twice this
from $10/3$ gives the first formula in (31). Substitute
$|K_t|=(1-t^3)/6$ and cancel $1-t$ to obtain the second.
Equations (30), (32), and (13) prove the bridge obstructions by taking
$t\downarrow0$.
\end{proof}

\section{Literature check, limitations, and verification}

\noindent\textbf{Rogers--Shephard and Zhang.}
B\"or\"oczky's 2005 paper proves linear Banach--Mazur stability for the
normalized Rogers--Shephard deficit, and a power $1/d$ estimate for the
polar projection (Zhang) deficit. Its Example 19 already uses truncation at
all simplex vertices and gives an obstruction beyond $1/(d-1)$ for that
different invariant. Truncation as a sharpness method is therefore prior
work. Equations (4)--(10) calculate the entry005 invariant and all its
maximum-simplex centroid distances specifically. Proposition \ref{prop:rs} shows why
the linear Rogers--Shephard result alone cannot give a linear bound in
$\ee$ in dimension three.

\medskip\noindent\textbf{Random simplices.}
Ambrus--B\"or\"oczky's random-simplex stability paper studies uniform
interior sampling; its stability theorem for the simplex as a maximizer is planar.
Entry005 instead uses determinants for the centered cone-volume law on
the polar boundary. No identity transporting the relevant defects between
these sampling models has been established here. The cited theorem is
therefore not a quantitative solution for this invariant in $d\ge3$.

\medskip\noindent\textbf{Inverse Minkowski.}
B\"or\"oczky--Machado--Ramos, arXiv:2603.17726v3, prove the sharp
Hausdorff exponent $1/(d-1)$ for inverse Minkowski stability in the
dual-convex distance of suitably dispersed normalized surface-area measures.
This is a plausible conversion step once closeness to a simplex measure
has been proved. It does not itself supply a bound on that measure distance
from the scalar gap $\ee$. This missing implication remains a separate
obligation, and no unconditional improved exponent is asserted from it.

\medskip\noindent\textbf{Minkowski asymmetry.}
Brandenberg, Gonz\'alez Merino, and Grundbacher (arXiv:2607.27041v2)
prove dimension-independent linear-order simplex approximation when
$d-s(K)$ is small. Equation (13) identifies that parameter exactly in
our family; (6) shows that a linear control of it by $\ee$ is impossible
in $d\ge3$. Thus this useful geometric theorem also requires a root-type
bridge in the present invariant.

\medskip\noindent\textbf{What is and is not proved.}
The exponent ceiling, exact invariant, maximum-simplex classification,
centroid excesses, translated excesses for those simplices, asymmetry,
and (31) are proved above. The general upper bound with exponent
$1/(d-1)$, the exact unrestricted distance in (11), and optimal universal
constants are not proved. No first-discovery or best-known claim is made.
The exact obstruction is compatible with the quantitative supplement's much smaller
positive exponent.

\medskip\noindent\textbf{Exact supporting computation.}
The accompanying standard-library script \texttt{check\_exact.py} uses
\texttt{fractions.Fraction} throughout. It checks facet determinants,
all vertex-subset simplex volumes, barycentric containment for several
different maximum simplices, and the dimension-three difference-body
volume by independently enumerating supporting facets and triangulating
them with rational arithmetic. These are finite regressions; the
all-dimensional statements and asymptotic conclusions rest on the written
proof, not on those computations. Errors are raised explicitly, so checks
remain active under \texttt{python3 -O}.

\begin{thebibliography}{9}
\bibitem{entry}
mxym, entry005, public versions 3--5, pinned commit
\nolinkurl{31e3d8a37e4a3051a7f5a2535be1c75642e15bcc}.
\url{https://github.com/mxym/math/tree/31e3d8a37e4a3051a7f5a2535be1c75642e15bcc/preprints/005-simplex-product-optimum}.
\bibitem{supplement}
mxym, \emph{Effective simplex rigidity for the projection cone invariant},
additive supplement to entry005, public source at commit
\nolinkurl{6785c1c830f8e19e2eb07b0bb89f4d475a8b154a}.
\url{https://github.com/mxym/math/blob/6785c1c830f8e19e2eb07b0bb89f4d475a8b154a/notes/quantitative-projection-simplex-stability/paper.md}.
\bibitem{rs}
K. B\"or\"oczky, Jr., \emph{The stability of the Rogers--Shephard inequality
and of some related inequalities}, Advances in Mathematics 190 (2005),
47--76. Author-hosted manuscript:
\url{https://www.renyi.hu/~carlos/rogerstab.pdf}.
\bibitem{random}
G. Ambrus and K. J. B\"or\"oczky, \emph{Stability results for the volume
of random simplices}, American Journal of Mathematics 136 (2014), 833--857.
\url{https://arxiv.org/abs/1005.5024}.
\bibitem{minkowski}
K. B\"or\"oczky, J. M. Machado and J. P. G. Ramos,
\emph{Quantitative Stability for Minkowski's problem}, arXiv:2603.17726v3,
13 May 2026. \url{https://arxiv.org/abs/2603.17726v3}.
\bibitem{asymmetry}
R. Brandenberg, B. Gonz\'alez Merino and F. Grundbacher,
\emph{Tight Stability Estimates Near the Simplex and Applications for
the Banach--Mazur Distance and Rogers--Shephard-Type Inequalities},
arXiv:2607.27041v2, 14 September 2026.
\url{https://arxiv.org/abs/2607.27041v2}.
\end{thebibliography}
\end{document}
